\src{174}
\begin{enumerate}
\item[\textup{i.}] $\varphi_{x_1,\ldots,x_p}$ depends in an alternating
$A$-multilinear fashion on the system of $x_i\in J$.
\item[\textup{ii.}] $\varphi_{x_1,\ldots,x_p}$ is zero whenever one of the
$x_i$ lies in $J^{2}$.
\end{enumerate}
In fact, ii follows from i, since $a\varphi_{\underline{x}}=0$ for $a\in J$, as
is seen by observing that all the modules in (3.2) are annihilated by $J$.

To define $\varphi_{\underline{x}}$, consider the complex
$K_{\underline{x}}$ whose underlying $A$-module is $\bigwedge A^{p}$ and whose
differential is contraction $i_{\underline{x}}$ by $\underline{x}$, regarded
as the linear form on $A^{p}$ with components $x_1,\ldots,x_p$. The differential
has degree $-1$, the degrees are nonnegative, and the homology of this complex
in dimension $0$ is $A/(x_1A+\cdots+x_pA)$. Since the $x_i$ belong to $J$, this
gives an augmentation $K_{\underline{x},0}\longrightarrow A/J$. Thus
$K_{\underline{x}}$ appears as an augmented free complex with augmentation
module $A/J$. In the familiar way, one obtains homomorphisms
\[
\operatorname{Ext}_{A}^{*}\bigl(H_0(K_{\underline{x}}),M\bigr)
\longrightarrow
H^{*}\bigl(\operatorname{Hom}_{A}(K_{\underline{x}},M)\bigr)
\]
and
\[
\operatorname{Tor}_{*}^{A}\bigl(H_0(K_{\underline{x}}),M\bigr)
\longleftarrow
H_{*}\bigl(K_{\underline{x}}\otimes_A M\bigr).
\]
Composing these with the homomorphisms on $\operatorname{Ext}$ and
$\operatorname{Tor}$ induced by the augmentation homomorphism
$H_0(K_{\underline{x}})\longrightarrow A/J$ gives homomorphisms
\begin{equation}\tag{3.3}
\left\{
\begin{aligned}
\psi_{\underline{x}}:
\operatorname{Ext}_{A}^{*}(A/J,M)
&\longrightarrow H^{*}\bigl(\operatorname{Hom}_{A}(K_{\underline{x}},M)\bigr),\\
\psi_{\underline{x}}:
\operatorname{Tor}_{*}^{A}(A/J,M)
&\longleftarrow H_{*}\bigl(K_{\underline{x}}\otimes_A M\bigr).
\end{aligned}
\right.
\end{equation}
One sees immediately that in the top dimension $p$, the homology of the
right-hand sides is $M/(x_1M+\cdots+x_pM)$ (respectively, the set of elements of
$M$ annihilated by every $x_i$). Since the $x_i$ lie in $J$, one obtains
homomorphisms
\begin{equation}\tag{3.4}
\left\{
\begin{aligned}
H^{p}\bigl(\operatorname{Hom}_{A}(K_{\underline{x}},M)\bigr)
&\longrightarrow M\otimes A/J,\\
H_{p}\bigl(K_{\underline{x}}\otimes_A M\bigr)
&\longleftarrow\operatorname{Hom}_{A}(A/J,M).
\end{aligned}
\right.
\end{equation}
Composing the homomorphisms (3.3) and (3.4) gives the homomorphisms (3.2) that we
wished to define. Verifying i is tedious but presents no difficulty.

